%! Project: Design and Conduct a Survey — Unit 1, Lesson 1.8 — Group Activity
\documentclass[10pt]{article}
\usepackage{saar-article}
\usepackage{saar-boxes}
\newcommand{\TallMath}[1]{$\displaystyle #1\rule[-1.4em]{0pt}{3.2em}$}

\begin{document}

\pageheader{Unit 1, Lesson 1.8}{Group Activity}
% NAMESTRIP: no \namedateperiod here. The student writes their name once, on the
% cover this component is stapled behind. See references/conventions.md.

\vspace{0.15in}

\begin{headlinebox}{frost}
\small
\textbf{Now you run one.} Your class picks \emph{one} question today and every
group asks the same question of every student in this room. Write your class
question here, exactly as the class agreed to it:

\vspace{4pt}
\writeline

\vspace{2pt}
The rest of this page is the study. Tier R plans it, Tier A builds the
instrument and collects, Tier E reports it and says what it is worth. The
Student Council's stratified survey of $60$ is on the board --- you will
compare against it at the end.
\end{headlinebox}

\vspace{0.03in}

\boxguard[26]
\begin{tcolorbox}[colback=white, colframe=black!40, breakable,
    title=\textbf{Tier R --- Remediate: Plan the Study},
    fonttitle=\bfseries, arc=2mm, left=3mm, right=3mm, top=1mm, bottom=1mm]
\small
\begin{enumerate}[leftmargin=1.6em, itemsep=3pt, topsep=3pt]

  \item Who is your \textbf{population} --- the group your answer will be about?
        \blank{6.4cm}

  \item How many people are in your \textbf{sample} (this room)? \blank{2cm}

  \item Fill in the plan, one stage at a time.

        \begin{center}
        \renewcommand{\arraystretch}{1.4}
        \begin{tabularx}{0.95\linewidth}{@{} p{2.5cm} X @{}}
        \toprule
        \textbf{Stage} & \textbf{Your class study} \\
        \midrule
        Formulate & \blank{7.2cm} \\
        Collect & \blank{7.2cm} \\
        Represent & \blank{7.2cm} \\
        Analyze \& report & \blank{7.2cm} \\
        \bottomrule
        \end{tabularx}
        \end{center}

  \item Which \textbf{collection method} are you using, and why that one?
        \writeline

\end{enumerate}
\end{tcolorbox}

\vspace{0.03in}

\boxguard[34]
\begin{tcolorbox}[colback=white, colframe=black!40, breakable,
    title=\textbf{Tier A --- Approaching Proficiency: Build It and Run It},
    fonttitle=\bfseries, arc=2mm, left=3mm, right=3mm, top=1.5mm, bottom=1.5mm]
\small
\begin{enumerate}[leftmargin=1.6em, itemsep=3pt, topsep=3pt]

  \item Write your \textbf{instrument} --- the item read word for word to every
        person, then its four choices. Check it against the four rules before
        you ask anyone.

        Item: \blank{11.2cm}

        Choices: (A) \blank{2.6cm} (B) \blank{2.6cm} (C) \blank{2.6cm}
        (D) \blank{2.6cm}

  \item \textbf{Collect.} Ask every student in the room and tally the answers.

        \begin{center}
        \renewcommand{\arraystretch}{1.5}
        \begin{tabularx}{0.95\linewidth}{@{} p{2.9cm} X c c @{}}
        \toprule
        \textbf{Choice} & \textbf{Tally} & \textbf{Count} & \textbf{Percent} \\
        \midrule
        (A) & & \blank{1.2cm} & \blank{1.6cm} \\
        (B) & & \blank{1.2cm} & \blank{1.6cm} \\
        (C) & & \blank{1.2cm} & \blank{1.6cm} \\
        (D) & & \blank{1.2cm} & \blank{1.6cm} \\
        \midrule
        \textbf{Total} & & \blank{1.2cm} & \blank{1.6cm} \\
        \bottomrule
        \end{tabularx}
        \end{center}

  \item Show the arithmetic for your \textbf{largest} category: its count
        divided by your total, as a percent.

        \begin{work}
          \dfrac{9}{25} &= 0.36 \\
                        &= 36\%
        \end{work}

  \item Show it for your \textbf{smallest} category.

        \begin{work}
          \dfrac{3}{25} &= 0.12 \\
                        &= 12\%
        \end{work}

  \item Do your four percents add to $100\%$? If not, find the mistake before
        you go on. \blank{5.2cm}

\end{enumerate}
\end{tcolorbox}

\vspace{0.03in}

\boxguard
\begin{tcolorbox}[colback=white, colframe=black!40, breakable,
    title=\textbf{Tier E --- Extension: Report It, Then Judge It},
    fonttitle=\bfseries, arc=2mm, left=3mm, right=3mm, top=1.5mm, bottom=1.5mm]
\small
\begin{enumerate}[leftmargin=1.6em, itemsep=4pt, topsep=3pt]

  \item If your class result held for all $900$ Riverbend students, about how
        many would choose your largest category?

        \begin{work}
          0.36 \times 900 &= 324 \text{ students}
        \end{work}

  \item Write the strongest \textbf{honest} sentence your study earns. Be
        careful about \emph{whom} it is allowed to describe.
        \writelines{2}

  \item The council's stratified sample of $60$ put club funding first at
        $35\%$. Your class of about $25$ put its own top choice first at a
        similar percent.

        \textbf{(a)} Whose result may be reported about all $900$ students?
        \blank{4.4cm}

        \textbf{(b)} The two studies nearly agree. Explain why that agreement
        still does not make your class's number trustworthy for the school.
        \writelines{2}

  \item Name one change that would let your class speak about all $900$
        students, and one thing it would cost. \writeline

\end{enumerate}
\end{tcolorbox}

\end{document}
